
\begin{exemples}
\label{XIII.4.4}
Obsérvese que, si~$S$ es un esquema conexo, las hipótesis
de~\Ref{XIII.4.1} se satisfacen cuando el morfismo~$f$ es propio y
plano, de presentación finita, con fibras geométricas separables y
conexas, siendo~$\LL$ arbitrario
(\cf~\Ref{X.1.3}). Las hipótesis de~\Ref{XIII.4.3} se
satisfacen si, además, $f$ es liso y posee una sección, y si se
designa por~$\LL$ el conjunto de los números primos distintos de las
características residuales de~$S$ (SGA~4 XV~2.1 y XVI~5.2).

Las hipótesis de~\Ref{XIII.4.1} también se satisfacen si~$S$ es conexo
y existe un esquema~$Z$ propio, de presentación finita y plano sobre~$S$,
con
\marginpar{420}
fibras geométricas separables y conexas, tal que~$X$ sea
el complementario en~$Z$ de un divisor de cruzamientos normales
relativamente a~$S$, siendo~$\LL$ el conjunto de los números
primos distintos de las características residuales de~$S$
\eqref{XIII.2.9}. Las hipótesis de~\Ref{XIII.4.3} se
satisfacen si, además, $f$ es liso y posee una sección.
\end{exemples}
